\section{An equality-level dictionary for the exterior operator}
\label{sec:operator-dictionary-v143}

The exterior contraction theorem is a supporting theorem for webs;
it is not an input to the principal pencil inverse theorem.  Its exact
relation to classical differential operators can nevertheless be
stated without a novelty slogan.  Fix the normalization
\[
 j_n(\epsilon\otimes\ell^n)=\bigwedge_{i=1}^n(\ell x_i),
 \qquad C_nj_n=2\,\mathrm{id},\qquad
 \kappa_{n,q}=\iota_qj_n.
\]
For \(h=x^{\mathsf t}Hx\) and \(q\) represented by \(Q\),
Theorem~\ref{thm:mixed-jacobian-contraction} identifies the actual
composite, on every homogeneous \(f\) of degree \(n\), as
\begin{equation}\label{eq:operator-dictionary-v143}
 C_n\varepsilon_h\kappa_{n,q}(\epsilon\otimes f)
 =\epsilon\otimes\left(
 \frac4n\sum_{ij}(HQ)_{ij}x_i\partial_j f
 -\frac{2h}{n(n-1)}\sum_{ij}Q_{ij}\partial_i\partial_j f
 \right).
\end{equation}
This formula compares operators with their coefficients and domains,
not merely isomorphic source representations.  It is valid for singular
\(q\) and arbitrary \(h\); it is not inferred from a nondegenerate
orthogonal change of variables.

For nondegenerate \(q\), choose \(h=\rho=q^{-1}\).  On degree
\(n\), the right hand side is exactly
\[
 \frac{2}{n(n-1)}\,
 \epsilon\otimes\bigl(E(E+n-2)-\rho\Delta_q\bigr)f.
\]
The operator in parentheses is the classical angular Casimir.
The harmonic decomposition and its eigenvalues are classical, not
additional representation-theoretic discoveries.  In the orthonormal
normalization \(q=I\), the induced symmetric-tensor norm is
\(\|x^\alpha\|^2=\alpha!/k!\) in degree \(k\).  The identity
\(j_n^*=2^{-n}C_n\), together with exterior adjunction, gives
\[
 \kappa_{n,q}^*\kappa_{n,q}
 =\frac{2^{1-n}}{n(n-1)}\bigl(E(E+n-2)-\rho\Delta_q\bigr).
\]
Thus the singular-value statement has a specified metric and a
specified scalar.  Changing from the unnormalized Fischer product
to symmetric-tensor norms without the factorial factors would not
be the same assertion.

For rank \(q=r<n\), choosing a radical square \(h=b^2\) in
\eqref{eq:operator-dictionary-v143} first forces \(\Delta_q f=0\)
for a vector in the kernel.  Varying \(h\) then forces dependence
only on radical variables.  This is the singular-rank argument of
Corollary~\ref{cor:singular-kernel-mixed}; no inverse quadratic form
or harmonic decomposition for a singular metric is used.

The closest inspected harmonic reference,
\cite[\S1, equations (1)--(3); Definition 6.1 and Lemma 6.2]{DeBieEelbodeRoels},
provides the Fischer decomposition and the multiplication--derivative
adjunction framework.  Its stated two-vector-variable operators act
on polynomial spaces; the exterior maps \(j_n,C_n,\iota_q\)
and the normalization in \eqref{eq:operator-dictionary-v143} must
still be identified.  The classical exterior-symmetric-square
summands in \cite[\S2, Theorem 3.5]{ChengWang} identify the source
module but do not fix that scalar normalization.  These are positive,
statement-specific comparisons with the cited sources.  They do not
establish that no other historical source contains the mixed identity.
All-rank kernel and support conclusions remain proved in full; their
geometric role is confined to the web recombination branch.
